%% ============================================================
\section{Limitations}
\label{sec:limitations}

SqC has several architectural limitations that cap its achievable
precision:

\textbf{No general preprocessor expansion.}  Macros still appear as
identifier tokens or function calls by default, and SqC does not
expand an arbitrary macro body.  A narrower, registry-based
macro-expansion engine now closes this gap for specific, recognized
idioms: it identifies function-like macros whose body frees and
null-writes their argument (the ``safe-free'' idiom, e.g.\
\texttt{SAFE\_FREE(p)}, feeding MEM30-C/MEM31-C) and macros whose body
writes to an output parameter (feeding EXP33-C), independent of the
macro's name.  Outside these two recognized shapes, an arbitrary
\texttt{SAFE\_FREE(p)}-style macro is still opaque.

We initially expected the remaining macro gap to be the largest
available precision lever, and report the negative result because it
redirects the obvious next step.  Feeding SqC the missing preprocessor
state directly---ingesting a \texttt{compile\_commands.json} to supply
the include search paths and \texttt{-D} definitions the source tree
never states---was implemented and measured, and it moved almost
nothing: $-0.09\%$ of findings from compile-database ingestion and
$-0.015\%$ from adding system includes, with no measurable movement in
either precision or recall.  The engine was never the bottleneck, and
neither was knowing where the headers live.  We therefore treat
general macro expansion as closed rather than as future work: the
residual FP mass sits elsewhere (caller contracts and ownership,
below), and a full preprocessing pass would buy a 10--100$\times$ TU
size increase and a per-rule regression surface for a fraction of a
percent.

\textbf{No general alias analysis.}  Pointer aliasing is largely
unresolved.  If \texttt{p} and \texttt{q} point to the same memory and
\texttt{q} is checked for null, SqC does not propagate that
information to \texttt{p} in the general case.  A field-sensitive
points-to module now tracks a narrower form of this---resolving
\texttt{\&var} and struct-field aliases for a handful of rules
(MEM30-C, MEM31-C, DCL13-C, ARR30-C)---but general pointer aliasing
remains unmodeled.

\textbf{No symbolic execution.}  SqC's value-range analysis handles
constant expressions but cannot reason about symbolic relationships
(e.g., that \texttt{n} is always less than \texttt{sizeof(buf)}).

\textbf{No SSA form.}  Without static single assignment, use-def
chains are approximate, limiting the precision of dataflow analyses.

\textbf{No ownership model.}  Cross-function memory ownership is only
partially tracked.  MEM31-C credits same-function ownership transfer
(allocate, then pass to a custom deallocator) via the points-to module
above, but does not distinguish a parameter-owned struct field
(populated by a function that does not own the struct's lifetime) from
a locally-owned one---the dominant remaining source of MEM31-C
real-world false positives (\cref{sec:realworld-results}).

These gaps are fundamental to the tree-sitter approach, though not
uniformly so: the macro and alias gaps above have narrowed via
targeted, rule-scoped analyses rather than general resolution.  A tool
built on Clang's AST would resolve macros and includes automatically
but would require a compilation database.  The current
architecture---AST plus CFG plus null-state dataflow plus value-range
analysis plus inter-procedural summaries (including cross-function
taint propagation)---has reached 87.1\% TP rate, exceeding the
75--80\% range this paper previously projected as achievable only with
alias analysis or symbolic execution.  In practice, most of that gain
came not from either of those two capabilities but from the opt-in
provenance-gated redesign of the INT rules
(\cref{sec:provenance-frontier}), which reframed the problem instead of
requiring deeper program analysis.  Further improvements past 87.1\%
likely still require some combination of general alias analysis,
symbolic execution, or a real ownership model, since the residual gaps
(CWE-369, CWE-121, the MEM31-C parameter-ownership case above) are
precisely the cases those three capabilities would address.

\textbf{Detection breadth, not output cleanliness, is the binding
constraint.}  The 87.1\% TP rate describes what fraction of SqC's
Juliet findings are real, not what fraction of Juliet's planted
defects SqC locates.  The latter---the flaw-hit rate, the share of
the suite's 132,406 flaw lines SqC actually flags---is 12.9\%, and it
has not moved across recent releases even as the TP rate and
real-world precision both improved.  Per-file detection (38.0\% of
flawed files) sits between the two: SqC flags something in over a
third of flawed files but lands on the specific planted line in about
an eighth of cases.  Reporting only the TP rate would materially
overstate the tool, and the flaw-hit rate is the honest headroom
signal.

The 13 scanned CWEs with zero detection (\cref{tab:precision-tiers})
reveal residual coverage gaps.  CWE-259 (Hard-Coded Password) is
blocked by tree-sitter's inability to expand preprocessor macros in
Juliet's Windows headers; CWE-328 (Reversible One-Way Hash) requires
a new cryptographic algorithm detection rule; the remainder (CWE-23,
CWE-123, CWE-176, CWE-321, CWE-366, CWE-570, CWE-571, CWE-667,
CWE-672, CWE-676, CWE-762) await dedicated rule implementations.

Finally, Juliet precision does not transfer to real-world precision.
To quantify the gap, we manually adjudicated a seeded random sample of
50 findings from each of the four noisiest rules in the v0.4.22
real-world run (MEM30-C, DCL13-C, INT32-C, EXP34-C; together 41K of
133K total findings).  Measured precision was 0\%, 34\%, 6\%, and 2\%
respectively---10.5\% combined---on rules whose Juliet results had
motivated no further work.  The FP modes are real-world idioms Juliet
rarely contains: free-then-reassign sequences, callback signatures
fixed by function-pointer contracts, loop counters structurally
bounded by data-structure sizes, and allocation checks placed on the
line after the allocation.  This is the strongest evidence for
treating Juliet as a rough approximation: a rule can be simultaneously
well-tuned on Juliet and produce thousands of false positives per
project in deployment.  A subsequent \emph{complete} audit of one
project---every INT32-C finding across all 220 SQLite files
adjudicated, rather than a 50-finding sample---sharpens the figure for
the worst rule to 0.3\% precision and isolates a single root cause,
operand provenance; we treat that case in detail in
\cref{sec:provenance-frontier}.

\subsection{The real-world corpus is the opposite population from the
intended one}
\label{sec:corpus-maturity}

The nine real-world projects are mature, actively maintained,
warning-clean C.  That makes them a demanding false-positive test, and
it is why findings from them have become upstream fixes.  It also makes
them a poor proxy for the population SqC is designed for, and this
biases every aggregate in a direction worth stating explicitly.

SqC requires no build system and tolerates source that does not
compile (\cref{sec:architecture}).  That is its distinguishing
property, and it implies that its nominal deployment is newer, less
mature, in-progress code wired into CI/CD early---not released
software.  Locating real defects in SQLite, curl and hostap
demonstrates the tool's reach, and we present it as exactly that: a
demonstration, not the representative case.

The consequence is that a per-rule real-world precision of 0\% is
frequently a fact about the corpus rather than about the rule.  Any
rule whose defect cannot survive review in a codebase of this maturity
is \emph{structurally} incapable of producing a true positive there.
DCL31-C is the clean example: 364 findings, 324 adjudicated, zero true
positives, which presents as 0.0\% precision.  That number measures
SqC's header reachability, not the rule's quality---mosquitto alone
falls from 1,365 DCL31-C findings with no \texttt{-I} to zero with
\texttt{-I /usr/include}.  The underlying defect is genuine: under C89
an implicit declaration causes the compiler to assume \texttt{int
f()}, so the return type is misread, argument checking does not
happen, and a returned pointer is truncated under LP64.  C99 removed
implicit declarations and C23 makes them an error.  The rule guards
real code; this corpus simply does not contain any.  Treating that
0.0\% as a rule-quality measurement is a category error.

Relatedly, rule \emph{applicability} is by design the user's decision,
expressed through manifest scoping and suppression, not something
detection logic adjudicates.  A Linux-only corpus should not exercise
the \texttt{WIN*} rules, and a zero there is correct rather than a
defect.

\textbf{How much this leaves unmeasured.}  Of 311 implemented rules,
186 have true-positive evidence from some benchmark---127 from Juliet,
144 from real-world TP/FN labels.  The remaining \textbf{125 (40\%)
have no true-positive evidence anywhere}: 65 fire on the corpus but
have only ever produced false positives (DCL31-C's class), and 60
never fire on the nine projects at all (including the categorically
inapplicable \texttt{WIN02-C} and \texttt{WIN30-C}).  Neither
benchmark can validate those rules, for the structural reason above,
so no amount of further adjudication on this corpus will resolve them.

The material for a third benchmark tier that could already exists in
the repository: 1,959 must-detect fixtures and 1,568 must-not-detect
fixtures across 309 rules, labeled by construction, of which 121 of
the 125 unvalidated rules already have a must-detect case.  Those
fixtures are currently consumed only as pass/fail unit tests and
contribute to no measured metric, so a rule can be thoroughly
exercised by tests and still present as having no detection evidence.
Scoring them---a miss in \texttt{fail/} is a false negative, a hit in
\texttt{pass/} a false positive---would measure detection where the
defect is guaranteed present, at no acquisition or adjudication cost.
We flag the gap rather than claim the result: that tier is not yet
built, and until it is, the honest statement about 40\% of SqC's rule
suite is that its detection capability is untested, not that it is
absent.
